% !TEX root = ../journal_0421052099.tex
% ---------------------------------------------------------------------
% Compresses the thesis environment chapter and the experimental-setup
% part of the evaluation chapter to ~1,600 words.
%
% TWO RULES THIS SECTION MUST NOT BREAK:
%
% 1. "pre-specified", NEVER "pre-registered". The protocol was fixed and
%    documented before the test sets were built; nothing was filed with a
%    registry.
%
% 2. Reachability is the SAMPLE ceiling, 97.0 / 99.2, not the full-split
%    97.3 / 99.66. Different populations. scripts/check_paper_numbers.py
%    fails the build if the full-split pair appears.
%
% September 2026 reframe: "certified" is defined where it is first used;
% the agentic baseline's depth, width, candidate caps and per-depth call
% count are stated, and the caps are attributed to THIS reimplementation
% rather than to the original paper; latency is named as a recorded
% measure with its cache caveat.
%
% CUT: the backbone selection reversal, the wrong-split incident, the
% API spend. One sentence points at the repository.
% ---------------------------------------------------------------------

\section{Experimental Setup}\label{sec:setup}

Everything in this section exists to make one inference safe: that a
difference between two systems is a difference between architectures. We
describe the environment all five systems share first, then the
benchmarks, the four baselines, and the single frozen backbone every one
of them calls, and finally the metrics and the protocol fixed before any
test question was run. Where something is \emph{not} held constant, we
name it here rather than leave it for a reader to find.

\subsection{Knowledge Environment}\label{sec:environment}

Every dataset this paper consumes is public, standard, and third-party.
The Freebase~\cite{bollacker2008freebase} snapshot and the two question
sets all arrive through one distribution: the per-question subgraphs
published with Reasoning-on-Graphs~\cite{luo2024reasoning}, released as
\textsf{rmanluo/RoG-webqsp} and \textsf{rmanluo/RoG-cwq}. The datasets
are used unmodified. What is built here is the environment derived from
them: a property graph of $2.59$ million entities and $8.31$ million
triples over $7{,}058$ distinct relation types, assembled as the union
of every per-question subgraph across both datasets and all three
splits. All five systems navigate that one graph. Just under two-thirds
of its nodes ($63.7\%$) are unnamed mediators encoding $n$-ary facts.
That is why the tool layer traverses them transparently rather than
exposing them as candidate entities (\Cref{sec:framework}).

An environment assembled this way is friendlier than full Freebase, and
we state the consequence rather than leave it to be inferred. Before any
system ran, a validation gate measured what fraction of questions have
at least one gold answer reachable from a topic entity within the
four-hop cap: $97.0\%$ on the WebQSP sample and $99.2\%$ on the
ComplexWebQuestions sample. A sample that passed that gate is what the
paper calls \emph{certified}. That ceiling is why a failure in
\Cref{sec:erroranalysis} can be attributed to the system rather than to
a missing edge. Almost every question is answerable in principle, so a
wrong answer is a navigation or selection error, not an absent fact. It
is equally why \textbf{absolute accuracies here do not transfer to a
full-Freebase setting}. Every comparison in \Cref{sec:results} is
between systems traversing the identical environment, which is what the
comparison is for. None of the numbers should be read against published
figures obtained over a different substrate; \Cref{sec:rog} shows what
such a reading looks like when it is done with the caveats attached.

One limit of the environment concerns expressiveness rather than
coverage. Every node carries a single entity label. There is no typed
literal, so a date or a number exists only incidentally, as an entity
whose name happens to be a numeric string, and the tool API exposes no
ordering or aggregation operator. A question whose gold answer is a
date, a measured quantity, or a superlative over candidates is therefore
unanswerable in this environment by construction.
\Cref{sec:erroranalysis} reports how much of the residue that class
accounts for, and \Cref{sec:discussion} carries it as a scope limit.

\subsection{Benchmarks, Baselines, and Backbone}\label{sec:baselines}

We evaluate on WebQSP~\cite{yih2016value} and
ComplexWebQuestions~\cite{talmor2018web}, the two standard multi-hop
KGQA benchmarks over Freebase. \textbf{We sample $400$ questions from
each} rather than running the full $1{,}628$ and $3{,}531$ test splits.
That was a pre-specified fallback, adopted when the full-split cost was
projected at roughly three times the budget available for the work. The
sample is stratified by hop count at a fixed seed, so its strata are
proportional to the split's: $256$, $128$, and $4$ questions at one,
two, and three or more hops on WebQSP, with $12$ whose gold answer the
gate found unreachable, and $137$, $211$, and $49$ on
ComplexWebQuestions, with $3$ unreachable. The resulting $95\%$
bootstrap intervals are roughly $\pm4$ to $\pm5$ points on Hits@1.
\Cref{sec:discussion} says which of the paper's claims survive that
width and which are reported as directions rather than effects. The two
benchmarks differ less in the typical question than in the tail. Exactly
half of the WebQSP sample carries a single gold answer, but its mean is
$7.15$ gold entities per question and its largest question lists $237$.
ComplexWebQuestions averages $1.75$ with a median of $1$. That tail is
one reason we report F1 beside Hits@1.

Four baselines span the design space. Each is a paradigm rather than a
tuned competitor, and three of the four are controls: they bound what a
paradigm can do on this environment, and nothing in this paper rests on
beating them. A \textbf{no-retrieval} control answers from parametric
memory alone, establishing what the backbone knows without the graph.
\textbf{Vector-RAG} retrieves verbalised triples by embedding
similarity. \textbf{Static GraphRAG}~\cite{edge2024local} fetches a
fixed-radius neighbourhood around the linked topic entities in one shot,
reranks it by embedding similarity to the question, and answers from the
top $30$ facts. Two properties of that implementation bound what it
licenses, and we state them here rather than leave a reader to infer
them from its scores. The radius is \emph{one logical hop}. Mediators
are traversed transparently, so a mediator path costs two graph edges
while remaining one hop, and a genuine two-hop chain is outside its
reach by construction. And expansion retrieves at most $100$ edges per
topic entity with no ordering criterion, so on a high-degree entity the
retained edges are an arbitrary sample rather than the best or the
first. At least one topic entity is truncated on $72.5\%$ of the $800$
test questions. This baseline therefore licenses one narrow claim:
\emph{retrieval fixed before reasoning must commit to a radius in
advance, and any fixed radius is wrong for some question in the set}. It
does not license a claim about what a two-hop static retriever would
score. That experiment was not run.

\textbf{Think-on-Graph}~\cite{sun2024think} is the agentic comparator
and the one comparison in this paper between two systems of the same
paradigm. It performs model-guided beam search over relations and
entities, and it is a reimplementation on this environment's tools
rather than a quoted number. Beam width and depth follow the original
paper, $3$ and $3$. At each depth it makes one relation-pruning call per
frontier entity, one entity-pruning call per selected relation, and one
sufficiency call, so a full depth at width $3$ costs up to thirteen
calls, and the $25$-call cap of \Cref{tab:budget} admits one full depth
and part of a second. Two widths beneath the beam are this
reimplementation's own rather than the original's: it offers the pruning
prompt the first $40$ relations of an entity and the first $20$
neighbours of a relation, against the $300$ and $200$ AGR sees
(\Cref{tab:budget}). Reimplementing is a stronger comparison than citing
the original figures, because graph, model, budget, and scoring rule are
held constant and only the algorithm varies. It also means its numbers
are not comparable to the published ones, and we do not present them as
such. All baselines were debugged on development data and certified
before any test question ran, against a pre-specified diagnostic
targeting the one failure that would be an implementation defect rather
than a paradigm limit: retrieval found the gold answer and the system
hedged anyway.

Every system shares one frozen backbone,
\textsf{gpt-5.4-mini-2026-03-17}, a dated snapshot rather than an alias,
at temperature $0.0$ with reasoning effort disabled. A check that logs
the reasoning-token count on every call read zero across every run in
this paper. The systems also share the entity resolver, the question
files, the output schema, and \emph{one budget meter with the same
$25$-call cap}. Both graph-navigating systems seed from the datasets'
annotated topic entities, which removes entity linking as a confound.
Holding all of this constant is what licenses the paper's central
inference: a difference between two systems is attributable to
architecture, not to model capacity, retrieval corpus, or spend. Three
things are deliberately not held constant, and we name them here rather
than bury them. The baselines receive a plain answering prompt, while
AGR's drafting prompt carries the grounding provisions that are part of
the system under test. The two navigators consider candidate sets of
different widths, the $40$ and $20$ above against $300$ and $200$. And
the shared call cap binds on the two navigators at very different rates,
which \Cref{sec:margin} shows is where the aggregate margin lives. The
second is where a reader should look first for a confound, and
\Cref{sec:margin} reports the measurement that bounds it.

\subsection{Metrics and Protocol}\label{sec:protocol}

Accuracy is reported as \textbf{Hits@1} and per-question \textbf{F1}.
Both systems and baselines return an unranked set of answer entities,
and Hits@1 here follows the benchmarks' convention for that case: a
question scores a hit if any returned entity is in the gold set. We
report F1 because a system can score a hit on that rule while asserting
several wrong entities alongside the right one, and \Cref{sec:margin}
checks directly whether AGR buys hits by naming more entities. A
\textbf{hedge} is an empty predicted entity set. Hedges score zero on
both metrics, and the hedge rate is reported as its own column, so a
system buying precision through abstention is visible rather than
flattered. Hits, hedges, and wrong answers always partition the whole
set. Nothing is dropped from the denominator, including questions the
environment provably cannot answer. Matching is exact after Unicode
normalisation, case folding, and whitespace stripping, with no fuzzy or
substring matching. We measured the obvious worry rather than assumed it
away: of AGR's wrong answers, $1$ of $65$ on WebQSP and $6$ of $99$ on
ComplexWebQuestions are surface-form near-misses that aggressive
normalisation would have scored as hits. The failure census of
\Cref{sec:erroranalysis} is reading real errors.

\textbf{Groundedness} is measured at two levels. The structural tier is
deterministic and applies to every asserted entity in all $4{,}000$
records: an entity is grounded if and only if it exists as a node in the
environment and connects to a topic entity within four hops. The
semantic tier asks the harder question of whether the entity is
supported \emph{as the answer}. For a stratified sample of $60$ answers
per system per dataset, it retrieves the graph's own shortest connecting
paths and has a language model judge whether any path expresses the
relationship the question asks about, under the governing instruction
that mere connectedness is not support. That judge was validated against
a human annotator on $100$ items, reaching $85\%$ agreement and Cohen's
$\kappa = 0.6995$. That \emph{misses} the $0.70$ we pre-specified, by
$0.0005$, and \Cref{sec:discussion} treats the miss as a limitation
rather than rounding it away. Because path retrieval takes only the
three shortest paths, semantic-tier rates are conservative lower bounds.
They are depressed uniformly across systems and read as a between-system
comparison rather than as absolute support rates. Efficiency is tokens
and model calls, both exact from the shared meter. Wall-clock time per
question was recorded as well, and \Cref{sec:cost} reports it with the
caveat it needs.

% \ref, not \Cref, for app:prompts: elsarticle's \appendix defines
% \thesection as "Appendix A", so \Cref printed "Section Appendix A".
% Plain \ref prints "Appendix A", which is the reading wanted.

The protocol was \textbf{pre-specified}. Research questions, metrics,
baseline certification thresholds, and ablation conditions were fixed
and documented before the test sets were constructed, and no
hyperparameter was tuned on test. Four decisions in particular were
written down before the measurements they govern: the frozen
configuration $\alpha = 0.7$ and $\tau = 0.20$ of \Cref{sec:framework},
selected on a disjoint $80$-question development set and never
revisited; the sample size of $400$ questions per dataset, named as the
fallback before the cost was known; the baseline certification
threshold; and the judge threshold of $\kappa \geq 0.7$. The last was
missed, and it is reported as missed. A threshold that counts only when
it is met is not a threshold. Nothing was filed with a public registry,
so the claim is pre-specification and not pre-registration. Confidence
intervals are $95\%$ percentile bootstrap intervals over questions from
$10{,}000$ resamples, and paired comparisons use the exact form of
McNemar's test on per-question correctness. The thirteen prompt
templates the five systems send, AGR's seven and the baselines' six, are
reproduced verbatim in \ref{app:prompts}, generated from the code by a
script so that they cannot drift. Deviations encountered during
execution, and there were several, are recorded in the repository rather
than narrated here.
